\section{Trabalhos Relacionados}
\label{sec:relacionados}

Estruturas do tipo DSU aparecem com frequência em problemas de conectividade incremental e em algoritmos clássicos como o de Kruskal para árvores geradoras mínimas. Nesse contexto, o trabalho de Tarjan estabelece a análise amortizada da combinação entre união por rank e compressão de caminho, obtendo o limite $\mathcal{O}(m\,\alpha(n))$ para uma sequência de $m$ operações \cite{tarjan1975efficiency}.

Em materiais didáticos e implementações voltadas a estudo, é comum comparar versões sem heurísticas com variantes mais otimizadas para evidenciar a diferença entre seus custos. O foco deste trabalho está nessa comparação, mas com um protocolo experimental explícito: as três variantes são avaliadas sobre as mesmas sequências reproduzíveis, com medição de tempo e de acessos à memória, para aproximar a análise assintótica do comportamento observado na execução.
